% Source: source/普通高考/2013/2013新课标2理(贵州,甘肃,青海,西藏,黑龙江,吉林,海南,宁夏,内蒙古,新疆,云南).pdf, page 1, question 6.
% Type: program flowchart; pdfimages -list reports no embedded bitmap (vector line art).
% Grid evidence: original page rendered at 180 dpi; grid step=100, fine_step=50.
% Node-edge list: 开始→输入N→k=1,S=0,T=1→T=T/k→S=S+T→k=k+1→k>N；
% k>N(是)→输出S→结束；k>N(否)→右侧回边→T=T/k入口主干汇合点。
% Solid segments: all node outlines and directed connectors; dashed segments: none.
% Labels 是/否 sit beside their corresponding decision exits. Every node uses
% local parindent=0pt and centered text. The loop arrow ends at the shared stem
% and does not create a second arrowhead on the downward main edge.
\begin{tikzpicture}[
  x=1cm,
  y=0.9cm,
  >=Stealth,
  line width=0.45pt,
  line cap=round,
  line join=round,
  every node/.style={font=\normalsize, inner sep=1pt, align=center, anchor=center,
    execute at begin node={\setlength{\parindent}{0pt}\noindent}},
  flow/.style={draw, line width=0.45pt, align=center, anchor=center,
    execute at begin node={\setlength{\parindent}{0pt}\noindent}},
  terminal/.style={flow, rounded corners=3pt, minimum width=1.25cm, minimum height=0.70cm},
  process/.style={flow, minimum height=0.72cm, align=center},
  io/.style={flow, trapezium, trapezium left angle=74, trapezium right angle=106,
    minimum height=0.74cm, align=center},
  decision/.style={flow, diamond, aspect=2.8, minimum width=3.35cm,
    minimum height=1.0cm, inner sep=1pt, align=center}
]
  \node[terminal] (start) at (0,10.55) {开始};
  \node[io, minimum width=2.25cm] (input) at (0,9.15) {输入 $N$};
  \node[process, minimum width=3.75cm] (init) at (0,7.75) {$k=1, S=0, T=1$};
  \node[process, minimum width=1.95cm, minimum height=1.05cm] (divide) at (0,6.20)
    {$T=\dfrac{T}{k}$};
  \node[process, minimum width=2.55cm] (add) at (0,4.75) {$S=S+T$};
  \node[process, minimum width=2.55cm] (increment) at (0,3.40) {$k=k+1$};
  \node[decision] (test) at (0,1.75) {$k>N$};
  \node[io, minimum width=2.25cm] (output) at (0,0.15) {输出 $S$};
  \node[terminal] (finish) at (0,-1.25) {结束};

  \coordinate (loopJoin) at (0,7.00);
  \coordinate (loopRight) at (2.65,1.75);

  % Main vertical path and terminating branch.
  \draw[->] (start.south) -- (input.north);
  \draw[->] (input.south) -- (init.north);
  \draw (init.south) -- (loopJoin);
  \draw[->] (loopJoin) -- (divide.north);
  \draw[->] (divide.south) -- (add.north);
  \draw[->] (add.south) -- (increment.north);
  \draw[->] (increment.south) -- (test.north);
  \draw[->] (test.south) -- node[left=2pt] {是} (output.north);
  \draw[->] (output.south) -- (finish.north);

  % 否 branch: right, upward, then a single left-pointing return arrow into
  % the main stem before T=T/k; the main stem retains its own downward arrow.
  \draw[->] (test.east) -- node[above=2pt] {否} (loopRight) -- (2.65,7.00) -- (loopJoin);
\end{tikzpicture}
